\documentclass[12pt]{article}
\usepackage{amsfonts,amssymb,amsmath,amsthm}
\usepackage[dvipsnames]{xcolor}
\usepackage[bottom=1in,top=1in,left=1in,right=1in]{geometry}
\usepackage[utf8]{inputenc}
\usepackage[english]{babel}
\usepackage{setspace}
\onehalfspacing
\usepackage[breaklinks=true]{hyperref}

\definecolor{darkblue}{rgb}{0,0,0.25}
\hypersetup{colorlinks=true,citecolor=blue,linkcolor=darkblue,urlcolor=darkblue,pdfpagemode=none}

\newtheorem{proposition}{Proposition}
\newtheorem{assumption}{Assumption}

\begin{document}

\section*{Appendix: Capital Wedge and the Induced Labor Wedge}

This appendix proves two claims in the weak two-sector environment with exogenous prices:
(i) manufacturing capital is inefficiently low in the decentralized equilibrium (DE) relative to the social planner (SP), and
(ii) this capital wedge implies a labor wedge.

%============================================================
% A. Environment
%============================================================

\subsection*{A. Environment}

There are two sectors: manufacturing ($m$) and services ($s$). Relative prices $(p_m/P,p_s/P)$ are treated as exogenous.

\paragraph{Services.}
Services are competitive and use labor only:
\[
Y_s = F_s(L_s), \qquad F_s'(L_s)>0,\quad F_s''(L_s)<0, \quad k_s\equiv 0.
\]
The service wage equals marginal product:
\begin{equation}\label{eq:ws}
w_s(L_s)=\frac{p_s}{P}F_s'(L_s),\qquad w_s'(L_s)=\frac{p_s}{P}F_s''(L_s)<0.
\end{equation}

\paragraph{Manufacturing.}
Manufacturing features search frictions and noncontractible capital. Per-match output is $f_m(k)$ with
\[
f_m'(k)>0,\qquad f_m''(k)<0.
\]
Let $\theta$ denote labor-market tightness and $q(\theta)$ the vacancy-filling rate, with
\[
q(\theta)>0,\qquad q'(\theta)<0,\qquad r+s_m>0.
\]
Let $z$ denote the unemployed flow payoff in manufacturing. Let $\beta\in(0,1)$ and impose Hosios $\beta=\eta$.

Define match surplus and the discount term:
\begin{equation}\label{eq:PiD}
\Pi_m(k)\equiv \frac{p_m}{P}f_m(k)-z,\qquad
D_m(\theta)\equiv r+s_m+(1-\beta)q(\theta)+\beta\,\theta q(\theta).
\end{equation}

%============================================================
% B. Equilibrium conditions and common locus
%============================================================

\subsection*{B. Manufacturing conditions and the common tightness locus}

The manufacturing DE satisfies the capital Euler equation and the free-entry (tightness) condition:
\begin{align}
P_k
&=
\frac{q(\theta_m)(1-\beta)}{r+s_m}\,
\frac{\frac{p_m}{P}f_m'(k_m)}{r+s_m+(1-\beta)q(\theta_m)},
\label{eq:DE_cap}\\
\frac{P_k k_m}{P}
&=
\frac{q(\theta_m)(1-\beta)}{r+s_m}\,
\frac{\Pi_m(k_m)}{D_m(\theta_m)}.
\label{eq:tight}
\end{align}
The planner's manufacturing capital Euler equation is
\begin{equation}\label{eq:SP_cap}
P_k
=
\frac{q(\theta_m)}{r+s_m}\,
\frac{\frac{p_m}{P}f_m'(k_m)}{r+s_m+q(\theta_m)}.
\end{equation}

Under CRS matching and free entry, Hosios implies the planner and DE share the same tightness condition \eqref{eq:tight}. Hence both allocations lie on the common tightness locus:
\begin{equation}\label{eq:common_locus}
\frac{P_k k_m}{P}
=
\frac{q(\theta_m)(1-\beta)}{r+s_m}\,
\frac{\Pi_m(k_m)}{D_m(\theta_m)}.
\end{equation}

%============================================================
% C. Sectoral indifference
%============================================================

\subsection*{C. Sectoral choice}

Manufacturing delivers the flow-equivalent return
\begin{equation}\label{eq:Rm}
R_m \equiv \beta\,\frac{\theta_m q(\theta_m)}{D_m(\theta_m)}\,\Pi_m(k_m).
\end{equation}
An interior allocation is characterized by sectoral indifference and labor feasibility:
\begin{equation}\label{eq:choice}
r(\epsilon_m-\epsilon_s)=R_m-w_s(L_s),\qquad L_m=\bar L-L_s.
\end{equation}

%============================================================
% D. Capital wedge
%============================================================

\subsection*{D. Capital wedge: $k_m^{DE}<k_m^{SP}$}

\begin{proposition}[Capital wedge]\label{prop:kwedge}
Assume $q(\theta)>0$, $q'(\theta)<0$, $f_m'(k)>0$, $f_m''(k)<0$, and impose Hosios $(\beta=\eta)$.
Define
\[
\Gamma^{DE}(\theta)\equiv (1-\beta)\frac{q(\theta)}{r+s_m+(1-\beta)q(\theta)},
\qquad
\Gamma^{SP}(\theta)\equiv \frac{q(\theta)}{r+s_m+q(\theta)}.
\]
Suppose that, along the relevant branch of the common locus \eqref{eq:common_locus}, the map
\[
k\mapsto \frac{\frac{p_m}{P}f_m'(k)}{r+s_m}\Gamma^j(\theta(k)),
\qquad j\in\{DE,SP\},
\]
is strictly decreasing in $k$.\footnote{A sufficient condition is that $\theta'(k)\ge 0$ along the common locus. More generally, the result only requires that any decline in $\theta(k)$ be sufficiently small that the concavity-induced decline in $f_m'(k)$ dominates the induced movement in the search term $\Gamma^j(\theta(k))$. Equivalently, it is enough that
\[
|f_m''(k)|\,\Gamma^j(\theta(k))
>
f_m'(k)\,|\Gamma_\theta^j(\theta(k))|\,|\theta'(k)|,
\qquad j\in\{DE,SP\}.
\]}
Then,
\[
k_m^{DE}<k_m^{SP}.
\]
\end{proposition}

\begin{proof}
The argument follows the first efficiency comparison in Acemoglu and Shimer (1999), adapted to the Hosios case.

Using \eqref{eq:DE_cap} and \eqref{eq:SP_cap}, the DE and planner capital Euler equations can be written as
\[
P_k=\frac{\frac{p_m}{P}f_m'(k)}{r+s_m}\Gamma^{DE}(\theta),
\qquad
P_k=\frac{\frac{p_m}{P}f_m'(k)}{r+s_m}\Gamma^{SP}(\theta).
\]
For any fixed $\theta$ and any $\beta\in(0,1)$,
\[
\Gamma^{DE}(\theta)<\Gamma^{SP}(\theta).
\]
Indeed, writing $q=q(\theta)>0$,
\[
(1-\beta)\frac{q}{r+s_m+(1-\beta)q}
<
\frac{q}{r+s_m+q}
\]
is equivalent to
\[
(1-\beta)(r+s_m+q)<r+s_m+(1-\beta)q,
\]
which reduces to
\[
(1-\beta)(r+s_m)<r+s_m,
\]
true since $\beta\in(0,1)$ and $r+s_m>0$.

Hence, holding $\theta$ fixed,
\[
\frac{\frac{p_m}{P}f_m'(k)}{r+s_m}\Gamma^{DE}(\theta)
<
\frac{\frac{p_m}{P}f_m'(k)}{r+s_m}\Gamma^{SP}(\theta),
\]
so the planner capital-Euler schedule lies strictly above the DE schedule at every common value of $\theta$.

Under Hosios, both allocations satisfy the same tightness condition \eqref{eq:common_locus}, so both are selected along the same common locus. By the maintained monotonicity condition, each capital-Euler right-hand side is strictly decreasing in $k$ along that locus near the interior solution. Since both schedules are set equal to the same constant $P_k$, and the planner schedule is pointwise above the DE schedule, the planner schedule reaches $P_k$ at a strictly larger capital stock. Therefore,
\[
k_m^{DE}<k_m^{SP}.
\]
\end{proof}

%============================================================
% E. Labor wedge
%============================================================

\subsection*{E. Labor wedge induced by the capital wedge}

\begin{assumption}[Monotonicity of manufacturing return on the common locus]\label{ass:returnmono}
Along the relevant branch of the common locus \eqref{eq:common_locus}, the product $k_m\theta_m$ is locally increasing in $k_m$; equivalently, if $\theta_m=\theta(k_m)$, then
\[
\theta(k_m)+k_m\theta'(k_m)>0
\]
locally around the interior solution.
\end{assumption}

\begin{proposition}[Labor wedge]\label{prop:Lwedge}
Suppose Assumption \ref{ass:returnmono} holds. Then, locally around an interior solution, $L_m$ is increasing in $k_m$. In particular,
\[
k_m^{DE}<k_m^{SP}\quad\Longrightarrow\quad L_m^{DE}<L_m^{SP}.
\]
\end{proposition}

\begin{proof}
At an interior allocation, sectoral indifference is given by
\begin{equation}\label{eq:choice_repeat}
r(\epsilon_m-\epsilon_s)=R_m-w_s(L_s),
\end{equation}
with labor feasibility $L_m=\bar L-L_s$. Differentiating \eqref{eq:choice_repeat} with respect to $k_m$ yields
\[
0=R_m'(k_m)-w_s'(L_s)\frac{dL_s}{dk_m},
\]
and hence
\begin{equation}\label{eq:dLsdk}
\frac{dL_s}{dk_m}=\frac{R_m'(k_m)}{w_s'(L_s)}.
\end{equation}
Because services are competitive and $F_s''(L_s)<0$, \eqref{eq:ws} implies
\[
w_s'(L_s)=\frac{p_s}{P}F_s''(L_s)<0.
\]
Therefore, the sign of $dL_s/dk_m$ is the opposite of the sign of $R_m'(k_m)$.

To sign $R_m'(k_m)$, use the common-locus condition \eqref{eq:common_locus}:
\[
\frac{P_k k_m}{P}
=
\frac{q(\theta_m)(1-\beta)}{r+s_m}\,
\frac{\Pi_m(k_m)}{D_m(\theta_m)}.
\]
Rearranging gives
\begin{equation}\label{eq:Pi_over_D_repeat}
\frac{\Pi_m(k_m)}{D_m(\theta_m)}
=
\frac{P_k k_m}{P}\cdot\frac{r+s_m}{(1-\beta)q(\theta_m)}.
\end{equation}
Substituting \eqref{eq:Pi_over_D_repeat} into the manufacturing return \eqref{eq:Rm},
\[
R_m
=
\beta\,\frac{\theta_m q(\theta_m)}{D_m(\theta_m)}\,\Pi_m(k_m),
\]
yields
\begin{equation}\label{eq:Rm_k_theta}
R_m
=
\beta\,\frac{P_k}{P}\cdot\frac{r+s_m}{1-\beta}\,k_m\theta_m.
\end{equation}
Thus, along the relevant branch of the common locus (where $\theta_m=\theta(k_m)$),
\begin{equation}\label{eq:Rprime_explicit}
R_m'(k_m)
=
\beta\,\frac{P_k}{P}\cdot\frac{r+s_m}{1-\beta}\,
\bigl[\theta(k_m)+k_m\theta'(k_m)\bigr].
\end{equation}
The multiplicative factor in \eqref{eq:Rprime_explicit} is strictly positive. By Assumption \ref{ass:returnmono},
\[
\theta(k_m)+k_m\theta'(k_m)>0,
\]
so \eqref{eq:Rprime_explicit} implies
\[
R_m'(k_m)>0.
\]

Returning to \eqref{eq:dLsdk}, since $R_m'(k_m)>0$ and $w_s'(L_s)<0$, we obtain
\[
\frac{dL_s}{dk_m}<0.
\]
Using labor feasibility, $L_m=\bar L-L_s$, it follows that
\[
\frac{dL_m}{dk_m}=-\frac{dL_s}{dk_m}>0.
\]
Hence $L_m$ is locally increasing in $k_m$.

Finally, Proposition \ref{prop:kwedge} established that $k_m^{DE}<k_m^{SP}$. Applying the local monotonicity of $L_m$ in $k_m$ yields
\[
L_m^{DE}<L_m^{SP}.
\]
\end{proof}
\end{document}